% ===================================================================
%  QUADRO N.º 4 — Idade Madura e Velhice.
%
%  Ceci n'est PAS une transcription : c'est une traduction, faite en
%  2026, en portugais d'aujourd'hui. Voir texto/pt/10-quadro-01.tex
%  pour la consigne, qui vaut pour les seize fichiers.
% ===================================================================

%%K t04-apar-1 apar
\VUsaut{4.0mm}
\VUtitre{15.0pt}{60}{QUADRO N.º 4}
\VUsaut{1.6mm}
\VUtitre{11.4pt}{0}{Idade Madura e Velhice.}
\VUsaut{3.2mm}

%%K t04-tit-1 sub
\VUcentre{11.4pt}{0}{\textit{Primeira cena.}}
\VUsaut{1.6mm}
\VUtitre{11.4pt}{0}{O Banquete de Casamento.}
\VUsaut{2.6mm}

%%K t04-tit-2 sub
\VUpk{10.2pt}{40}{O outono.}
\VUsaut{3.0mm}

%%K t04-c1-01-1 p
1. --- Os jovens cresceram. Um deles, João, casa-se. Assistimos ao
\VUgras{banquete de casamento}, que reúne à volta da mesma mesa as
famílias dos dois esposos.

%%K t04-c1-02-1 p
2. --- Estamos no \VUgras{outono}, no mês de \VUgras{setembro}. O vento frio
de \VUgras{outubro} e de \VUgras{novembro} ainda não despojou os campos da
sua verde roupagem. O ar da tarde é tépido e perfumado; por isso o jantar
é servido sob a \VUgras{galeria} \textsuperscript{(8)} que dá para o jardim.

%%K t04-c1-03-1 p
3. --- Ao longe, na \VUgras{colina} \textsuperscript{(3)}, está a casa dos pais de
João, uma elegante \VUgras{casa de campo} \textsuperscript{(1)} escondida entre a
verdura; belas vinhas, que dão um vinho excelente, cobrem a encosta do
outeiro. O vinhateiro cultiva-as e trata delas.

%%K t04-c1-04-1 p
4. --- Por cima das vinhas passa um bando de \VUgras{perdizes}
\textsuperscript{(2)}, assustadas pela chegada do \VUgras{guarda-caça}
\textsuperscript{(5)}. Este, com a \VUgras{espingarda} debaixo do braço e o
\VUgras{surrão a tiracolo}, bate os \VUgras{matagais} \textsuperscript{(4)} com o
seu \VUgras{cão} \textsuperscript{(6)}. Vigia a herdade e impede que os furtivos
acabem com a caça.

%%K t04-c1-05-1 p
5. --- Fora da \VUgras{galeria} trepa uma \VUgras{videira}, da qual pendem
pesados \VUgras{cachos de uva} \textsuperscript{(7)}.

%%K t04-c1-06-1 p
6. --- As bonitas flores de outono que enfeitam a \VUgras{mesa}
\textsuperscript{(16)} são \VUgras{crisântemos} \textsuperscript{(9)}; o \VUgras{pai}
\textsuperscript{(10)} da noiva pôs um na lapela do fraque.

%%K t04-c1-07-1 p
7. --- A refeição chega ao fim. O \VUgras{criado} \textsuperscript{(11)} começa a
trazer as sobremesas e vai retirar as diferentes peças do talher ---
\VUgras{colheres} \textsuperscript{(14)}, \VUgras{garfos} \textsuperscript{(12)},
\VUgras{facas} \textsuperscript{(13)}, \VUgras{travessas} \textsuperscript{(15)} --- que
já não fazem falta.

%%K t04-tit-3 sub
\VUpk{10.2pt}{40}{A Família.}
\VUsaut{3.0mm}

%%K t04-c1-08-1 p
8. --- Todos parecem alegres, e o sorriso com que a \VUgras{mãe}
\textsuperscript{(17)} do noivo olha para os seus filhos prova a sua
felicidade.

%%K t04-c1-09-1 p
9. --- Este casamento une duas famílias ricas da região. João, o
\VUgras{marido} \textsuperscript{(18)}, é \VUgras{filho} de um grande
proprietário e torna-se \VUgras{genro} de um grande industrial.

%%K t04-c1-10-1 p
10. --- Os \VUgras{esposos} são jovens e no entanto estavam
\VUgras{noivos} havia muito tempo. A \VUgras{jovem esposa}
\textsuperscript{(19)}, Marta, está encantadora com o seu \VUgras{vestido branco}
enfeitado com flores de laranjeira.

%%K t04-c1-11-1 p
11. --- O seu \VUgras{sogro} \textsuperscript{(20)} está muito contente de ter uma
\VUgras{nora} tão amável, e a \VUgras{sogra} \textsuperscript{(21)} de João sorri
ao ver a sua \VUgras{filha} tão bela.

%%K t04-c1-12-1 p
12. --- Na ponta da mesa sentam-se os demais \VUgras{convidados}
\textsuperscript{(22)}: \VUgras{parentes, amigos, primos e primas}. Um
\VUgras{mordomo} \textsuperscript{(23)}, com o guardanapo debaixo do braço,
serve-os com diligência.

%%K t04-tit-4 sub
\VUpk{10.2pt}{40}{Os Amigos.}
\VUsaut{3.0mm}

%%K t04-c1-13-1 p
13. --- Do lado direito da mesa está sentada uma \VUgras{menina} \textsuperscript{(24)},
\VUgras{amiga} da noiva, e depois o \VUgras{irmão} desta, que é o seu
\VUgras{padrinho de casamento} \textsuperscript{(25)}. Conversa com a sua
\VUgras{dama de honor} \textsuperscript{(26)}, irmã do noivo, que por este
casamento se torna \VUgras{cunhada} de Marta. É uma jovem encantadora.
Traz belas \VUgras{joias}, um \VUgras{colar de pérolas} e uma
\VUgras{pulseira} de ouro.

%%K t04-c1-14-1 p
14. --- Perto deles, o senhor que está de pé e levanta o seu copo é um
\VUgras{amigo} \textsuperscript{(27)} da família. É uma das \VUgras{testemunhas do
noivo}. Faz um \VUgras{brinde}, bebe à \VUgras{saúde} dos recém-casados e
deseja-lhes muita felicidade. Vai a rigor, de fraque e \VUgras{sapatos de
verniz} \textsuperscript{(28)}. Acompanha a \VUgras{avó} \textsuperscript{(29)}, que é
\VUgras{viúva}.

%%K t04-c1-15-1 p
15. --- O jovem de \VUgras{fraque} \textsuperscript{(31)}, do fino bigode preto, é
o \VUgras{cunhado} \textsuperscript{(30)} de João. É um engenheiro distinto. Está
ao lado de uma \VUgras{jovem senhora} \textsuperscript{(33)} muito elegante, a
\VUgras{irmã mais velha} de Marta e mãe da linda menina que vem, pelo
braço do seu primo, oferecer-lhe uma flor e \VUgras{dar-lhe} as
\VUgras{boas-noites}. Esta senhora traz uns \VUgras{brincos} muito belos e
um elegante \VUgras{penteado}.

%%K t04-c1-16-1 p
16. --- O priminho, tão gracioso com o seu \VUgras{blusão}
\textsuperscript{(36)} e as suas \VUgras{calças} \textsuperscript{(37)} abaloadas, dá
muita pena. É \VUgras{órfão} \textsuperscript{(35)}. Perdeu muito novo o pai e a
mãe. O seu tio é o seu \VUgras{tutor} \textsuperscript{(38)} e ficou solteiro para
criar o pobre menino. É um homem excelente. É um tanto calvo e muito
míope; usa \VUgras{pincenê}.

%%K t04-tit-5 sub
\VUsaut{2.4mm}
\VUpk{10.2pt}{40}{A Sobremesa.}
\VUsaut{3.0mm}

%%K t04-c1-17-1 p
17. --- Atrás dos comensais, numa mesa coberta com uma \VUgras{toalha},
está preparada a \VUgras{sobremesa} \textsuperscript{(39)}. Numa taça grande há
fruta: \VUgras{pêssegos} \textsuperscript{(40)}, \VUgras{peras}
\textsuperscript{(41)}, \VUgras{maçãs} \textsuperscript{(42)}, \VUgras{damascos}
\textsuperscript{(43)} e \VUgras{uvas} \textsuperscript{(44)}. Ao lado,
\VUgras{ananases} \textsuperscript{(45)}, um belo \VUgras{bolo}
\textsuperscript{(46)}, \VUgras{garrafas de licor} \textsuperscript{(48)} e
\VUgras{champanhe} \textsuperscript{(49)}, e também pilhas de \VUgras{pratos}
\textsuperscript{(47)} limpos.

%%K t04-c1-18-1 p
18. --- O \VUgras{escanção} \textsuperscript{(50)} desarrolha uma \VUgras{garrafa}
\textsuperscript{(52)} de vinho velho com um \VUgras{saca-rolhas}
\textsuperscript{(51)}. A \VUgras{rolha} \textsuperscript{(53)} parece muito difícil de
tirar. Este escanção traz um \VUgras{guardanapo} \textsuperscript{(54)} debaixo
do braço.

%%K t04-c1-19-1 p
19. --- Depois do jantar, os senhores passarão à sala de fumo e as
senhoras irão arranjar-se para o baile.

%%K t04-tit-6 sub
\VUsaut{5.0mm}
\VUcentre{11.4pt}{0}{\textit{Segunda cena.}}
\VUsaut{1.6mm}
\VUpk{11.4pt}{40}{O Aniversário do avô.}
\VUsaut{2.6mm}

%%K t04-tit-7 sub
\VUpk{10.2pt}{40}{A Visita.}
\VUsaut{3.0mm}

%%K t04-c2-01-1 p
1. --- Passaram dez anos. Os pais de João envelheceram e gostam muito
dos seus três netos, dois \VUgras{netos} \textsuperscript{(2)} e uma \VUgras{neta}
\textsuperscript{(3)}. Como hoje é o \VUgras{aniversário do avô}, os filhos vêm
felicitá-lo.

%%K t04-c2-02-1 p
2. --- O pequeno Pedro é ainda muito criança; tem apenas três anos. Vê
que se aproxima o \VUgras{gato} \textsuperscript{(1)}. Quer acariciar o seu belo
\VUgras{pelo} sedoso e a sua longa \VUgras{cauda}; não pensa que, sob essas
graciosas \VUgras{patas}, se escondem feias unhas pontiagudas que bem
poderiam arranhá-lo.

%%K t04-c2-03-1 p
3. --- A sua irmã Gabriela, que lhe dá a mão, é muito mais séria, e
Marcelo, com o seu grande \VUgras{casacão} \textsuperscript{(4)} e o seu
\VUgras{cinto} \textsuperscript{(5)} de couro, parece já um homenzinho, apesar
das calças curtas que deixam ver as suas \VUgras{meias}
\textsuperscript{(6)}.

%%K t04-c2-04-1 p
4. --- O pai deles pôs o seu grosso \VUgras{sobretudo} \textsuperscript{(7)} de
inverno com \VUgras{gola de pele} \textsuperscript{(8)} e, no \VUgras{bolso}
\textsuperscript{(9)}, leva um lenço. A mulher dele traz o seu chapéu de feltro
enfeitado com \VUgras{plumas} \textsuperscript{(10)}, o seu \VUgras{casaco}
\textsuperscript{(11)}, a sua \VUgras{boa de pele} \textsuperscript{(12)} e o seu
\VUgras{regalo} \textsuperscript{(13)}.

%%K t04-c2-05-1 p
5. --- Está muito frio, porque reina o inverno. A \VUgras{neve}
\textsuperscript{(19)} caiu todo o dia e os telhados das casas estão brancos.
Com a \VUgras{noite}, o frio torna-se ainda mais vivo. No céu, a
\VUgras{lua} \textsuperscript{(17)} e as \VUgras{estrelas} \textsuperscript{(18)} brilham
e atravessam a \VUgras{escuridão}.

%%K t04-tit-8 sub
\VUsaut{4.2mm}
\VUpk{10.2pt}{40}{A Doença.}
\VUsaut{3.0mm}

%%K t04-c2-06-1 p
6. --- O pobre \VUgras{cego} \textsuperscript{(20)} que atravessa a praça diante da
\VUgras{farmácia} \textsuperscript{(16)}, guiado pelo seu \VUgras{caniche}
\textsuperscript{(21)}, é muito desditoso. Justina, a \VUgras{criada}
\textsuperscript{(14)}, traz da cozinha uma \VUgras{tigela} \textsuperscript{(15)} de
\VUgras{chá} muito quente e bem açucarado.

%%K t04-c2-07-1 p
7. --- A \VUgras{avó} \textsuperscript{(23)}, que quer ir buscar os seus
\VUgras{óculos} \textsuperscript{(26)} à mesa para ler a \VUgras{receita}
\textsuperscript{(27)} que o \VUgras{médico} \textsuperscript{(25)} acaba de escrever,
levanta-se da sua poltrona apoiando-se na sua \VUgras{bengala}
\textsuperscript{(24)}.

%%K t04-c2-08-1 p
8. --- Sofre muitas vezes de pequenas \VUgras{indisposições},
\VUgras{tonturas}, \VUgras{constipações} e \VUgras{dores de dentes}; tem as
pernas fracas e a vista já não é muito boa. Apesar disso, troça de bom
grado do seu velho \VUgras{médico}, que sempre lhe quer receitar uma
\VUgras{poção} \textsuperscript{(28)}. Hoje está muito contente por ter a sua
jovem família à sua volta.

%%K t04-c2-09-1 p
9. --- Está-se bem no quarto bem fechado. Na \VUgras{lareira}
\textsuperscript{(29)} de mármore arde um \VUgras{fogo esplêndido}
\textsuperscript{(30)}, e sente-se uma pessoa feliz com a doce \VUgras{luz} do
\VUgras{candeeiro} \textsuperscript{(31)} que acabam de acender.

%%K t04-tit-9 sub
\VUsaut{4.2mm}
\VUpk{10.2pt}{40}{O Avô.}
\VUsaut{3.0mm}

%%K t04-c2-10-1 p
10. --- Até o \VUgras{avô} \textsuperscript{(33)} se esquece de que esteve doente,
de que o seu \VUgras{reumatismo} o prende à sua \VUgras{poltrona}
\textsuperscript{(34)} e de que, ainda ontem, tinha \VUgras{febre e desmaios}.
Já não se lembra de que no inverno passado teve uma \VUgras{pneumonia} e
de que uma \VUgras{tosse} rouca e penosa o fazia sofrer muito.

%%K t04-c2-10-2 p
E no entanto restabeleceu-se com muita dificuldade. Agora está um pouco
melhor. Mas a perna, que apoia sobre uma \VUgras{almofada}
\textsuperscript{(38)} posta num pequeno \VUgras{banco} \textsuperscript{(39)}, ainda
lhe dói.

%%K t04-c2-11-1 p
11. --- Quando está bem embrulhado no seu quente \VUgras{roupão}
\textsuperscript{(35)} de grossos \VUgras{botões} \textsuperscript{(36)} de madrepérola,
e quando tem a seu lado, para lhe ler o jornal, a pequena \VUgras{órfã}
\textsuperscript{(40)} com a sua \VUgras{touca} branca, o seu \VUgras{vestido}
\textsuperscript{(42)} e a sua \VUgras{esclavina} \textsuperscript{(41)} azuis, não se
queixa.

%%K t04-c2-12-1 p
12. --- Todos os anos, quando o \VUgras{calendário} \textsuperscript{(43)} traz de
novo a \VUgras{data} do primeiro de \VUgras{dezembro}, espera com
impaciência os seus \VUgras{netos}, porque sabe que não esquecerão essa
\VUgras{data} importante.

%%K t04-c2-13-1 p
13. --- Recorda com emoção os tempos longínquos em que ele próprio ia
felicitar o glorioso \VUgras{marechal} do Império, cujo \VUgras{retrato}
\textsuperscript{(44)} pende da parede, e que é o \VUgras{bisavô} das crianças.
\VUsaut{6.0mm}
\VUfilet{22mm}
